\documentclass[english, parskip=full]{scrartcl}
\usepackage[ngerman]{babel}  % Sprache auf Deutsch 
\usepackage[utf8]{inputenc}  % Kodierung der Datei
\usepackage[left=2cm, right=2cm, top=2cm, bottom=3cm, footskip=1.5cm]{geometry}
\input{myscrarticle}

\newcommand*{\titel}{Aufgaben}
\newcommand*{\autor}{Joachim Schwardt}

\hypersetup{pdfauthor={\autor}, pdftitle={\titel}}

%\titlehead{titlehead}
\subject{Geometrie}
\title{\titel}
\author{\autor}
\date{\today}

\begin{document}

%\begin{titlepage}
%\maketitle  % Titel setzen
%\tableofcontents  % Inhaltsverzeichnis setzen
%\end{titlepage}

\section{Dreiecke und Analysis}
Die erste Aufgabe befasst sich mit einem Konstrukt bestehend aus zwei Dreiecken, das in \autoref{figure:tikz skizze aufgabe dreieck alpha beta} skizziert ist.
Für beliebige Winkel $\alpha,\beta$ ist nach dem unbekannten Winkel $\epsilon$ gesucht.\\
\begin{figure}[!ht]
\centering
\begin{tikzpicture}   
    \pgfmathsetmacro{\a}{3.0}
    \pgfmathsetmacro{\aa}{45}
    \pgfmathsetmacro{\ba}{30} 
    \pgfmathsetmacro{\r}{0.05}
    
    
    \coordinate (o) at (0,0);
    \coordinate (a) at (\a,0);
    \coordinate (a2) at ($(a) + (180-\aa:\a)$);
    \coordinate (b) at (2*\a,0);
    \coordinate (b2) at ($(b) + (180-\ba:\a)$);
    \coordinate (c) at (-\a,0);
    \coordinate (e) at (intersection of a--a2 and b--b2);
    
    \foreach \v in {o,a,b,e} {
        \draw[fill] (\v) circle (\r);    
    }
    \draw[thick] (o) -- (a) node[below,midway] {1} -- (b) node[below,midway] {1}  -- (e) -- (o);    \draw[thick] (e) -- (a);
    \draw[thick] (e) -- (a);
    \pic["$\alpha$",draw=black,angle eccentricity=0.6,angle radius=1cm] {angle=e--a--o};
    \pic["$\beta$",draw=black,angle eccentricity=0.8,angle radius=1cm] {angle=e--b--a};
    \pic["$\epsilon$",draw=black,angle eccentricity=0.6,angle radius=1cm] {angle=o--e--a};    
\end{tikzpicture}
\caption{Skizze der Aufgabenstellung, hier für den Spezialfall $\alpha=\frac{\pi}{4}$ und $\beta = \frac{\pi}{6}$.
Man beachte, dass die unteren Kanten der Dreiecke dieselbe, und ohne EInschränkung als Eins definierte Länge haben.
Gesucht ist der Winkel $\epsilon$.}
\label{figure:tikz skizze aufgabe dreieck alpha beta}
\end{figure}
Für bestimmte Winkel ist eine konstruktive Lösung vermutlich eleganter, im Allgemeinen lassen sich Probleme dieser Art aber auch gut mit etwas Transfer aus Analysis bearbeiten.
Zunächst führen wir noch den Winkel $\gamma$ ein und erweitern die Skizze um ein paar hilfreiche Bezeichnungen, wie in \autoref{figure:tikz skizze loesung dreieck alpha beta} dargestellt.
\begin{figure}[!ht]
\centering
\begin{tikzpicture}   
    \pgfmathsetmacro{\a}{3.0}
    \pgfmathsetmacro{\aa}{45}
    \pgfmathsetmacro{\ba}{30} 
    \pgfmathsetmacro{\r}{0.05}
    
    
    \coordinate (o) at (0,0);
    \coordinate (a) at (\a,0);
    \coordinate (a2) at ($(a) + (180-\aa:\a)$);
    \coordinate (b) at (2*\a,0);
    \coordinate (b2) at ($(b) + (180-\ba:\a)$);
    \coordinate (c) at (-0.4*\a,0);
    \coordinate (e) at (intersection of a--a2 and b--b2);
    
    \foreach \v in {o,a,b,e} {
        \draw[fill] (\v) circle (\r);    
    }
    \draw[thick] (o) -- (a) node[below,midway] {1} -- (b) node[below,midway] {1}  -- (e) -- (o);    \draw[thick] (e) -- (a);
    \pic["$\alpha$",draw=black,angle eccentricity=0.6,angle radius=1cm] {angle=e--a--o};
    \pic["$\beta$",draw=black,angle eccentricity=0.8,angle radius=1cm] {angle=e--b--a};
    \pic["$\epsilon$",draw=black,angle eccentricity=0.6,angle radius=1cm] {angle=o--e--a};
    
    %solution
    \draw[thick,blue] (o) -- (c);
    \pic["$\gamma$",draw=blue,blue,angle eccentricity=0.6,angle radius=1cm] {angle=e--o--c};
    \node at (e) [above right,blue] {$E=(x,y)$};
    \node at (o) [below,blue] {$O$};
    \node at (a) [below,blue] {$A$};
    \node at (b) [below,blue] {$B$};
    
\end{tikzpicture}
\caption{Skizze der Aufgabenstellung erweitert mit für die Lösung hilfreichen Konstruktionen, hier für den Spezialfall $\alpha=\frac{\pi}{4}$ und $\beta = \frac{\pi}{6}$.
Gesucht ist der Winkel $\epsilon$.}
\label{figure:tikz skizze loesung dreieck alpha beta}
\end{figure}\\
Stellt man sich die Skizze um 90 Grad im Uhrzeigersinn gedreht vor, dann lassen sich die drei Geraden $\overline{OE}$, $\overline{AE}$ und $\overline{BE}$ durch die Funktionen
\begin{align}
f_1(t) &= ct, \qquad f_2(t) = at - 1 \qquad \text{und} \qquad f_3(t) = bt - 2
\end{align}
beschreiben.
Dabei gilt jeweils $f_j(x) = y$ für $j\in\klc*{1,2,3}$ und die Anstiege der Geraden entsprechen
\begin{align}
a &= \frac{1}{\tan(\alpha)}, \qquad b = \frac{1}{\tan(\beta)}\qquad \text{und}\qquad c = \frac{1}{\tan(\gamma)},
\end{align}
was sich beispielsweise an den Steigungsdreiecken ablesen lässt.
Zunächst bestimmen wir die Koordianten das Punktes $E$ durch das Lösen des Gleichungssystems
\begin{align}
f_2(x) &= ax - 1 \overset{!}{=} y,\label{eqn:gls fuer Punkt E f2}\\
f_3(x) &= bx - 2 \overset{!}{=} y\label{eqn:gls fuer Punkt E f3}.
\end{align}
Die Subtraktion von \autoref{eqn:gls fuer Punkt E f2} und \autoref{eqn:gls fuer Punkt E f3} ergibt zunächst
\begin{align}
(a-b)x +1 &= 0
\end{align}
und somit $x= \frac{1}{b - a}$.
Damit finden wir dann
\begin{align}
y &= ax - 1 = \frac{a}{b-a} - 1 = \frac{2a - b}{b - a}.
\end{align}
Man beachte nun, dass der gesuchte Winkel über die Innenwinkelsumme von $\Delta OAE$ zu
\begin{align}
\epsilon &= \pi - \alpha - (\pi - \gamma) = \gamma - \alpha
\end{align}
bestimmt werden kann.
Da $\alpha$ gegeben  ist benötigen wir jetzt also nur noch $\gamma = \arctan \kla*{\frac{1}{c}}$.
De Parameter $c$ erhalten wir dabei aus der Bedingung $f_1(x) = y$.
Einsetzen liefert
\begin{align}
c &= \frac{y}{x} = \frac{2a-b}{b-a} (b-a) = 2a-b.
\end{align}
Mit $\arctan(x) = \arccot\kla*{\frac{1}{x}}$ kann man das Endergebnis in kompakter -- aber, wie wir noch sehen werden, eher unhandlicher -- Form als
\begin{align}
\epsilon &= \arccot \kla*{2a-b} - \alpha = \arccot\kla*{2\cot(\alpha) - \cot(\beta)} - \alpha \label{eqn:ergebnis epsilon v1}
\end{align}
schreiben.
Für den Spezialfall $\alpha = \frac{\pi}{4}$ und $\beta = \frac{\pi}{6}$ lautet diese Formel allerdings
\begin{align}
\epsilon &= \arccot\kla*{2 - \sqrt{3}} - \frac{\pi}{4}
\end{align}
und der erste Term scheint auf den ersten Blick keine Vereinfachung zuzulassen. 
Tatsächlich kann man zeigen, dass $\arccot(2-\sqrt{3}) = \frac{5\pi}{12}$ gilt, einfacher ist allerdings eine Manipulation von \autoref{eqn:ergebnis epsilon v1}.
Wir bedinen uns dafür der Additionstheoreme für Sinus und Kosinus, die sich in jeder Formelsammlung finden lassen sollten.
Damit gilt zunächst allgemein
\begin{align}
\tan(x+y) &= \frac{\sin(x+y)}{\cos(x+y)} = \frac{\sin(x) \cos(y) + \cos(x) \sin(y)}{\cos(x) \cos(y) - \sin(x) \sin(y)} \color{blue} \frac{\frac{1}{\cos(x)\cos(y)}}{\frac{1}{\cos(x)\cos(y)}} \color{black} = \frac{\tan(x) + \tan(y)}{1 - \tan(x)\tan(y)}. \label{eqn:additionstheorem tan}
\end{align}
Damit und unter Verwendung von $\tan(x) = \frac{1}{\cot(x)}$ findet man dann
%\begin{align}
%\tan(\epsilon) &= \tan\kla*{\arctan\kla*{\frac{1}{2a - b}} - \alpha} = \frac{\frac{1}{2a - b} - \tan(\alpha)}{1 + \tan(\alpha) \frac{1}{2a - b}} = \frac{1 - \tan(\alpha) (2a - b)}{2a - b + \tan(\alpha)}.
%\end{align}
%Erweitern des Bruches mit $\frac{1}{ab} = \tan(\alpha) \tan(\beta)$ ergibt
%\begin{align}
%\tan(\epsilon) &= \frac{1 - \tan(\alpha) \kla*{\frac{2}{\tan(\alpha)} - \frac{1}{\tan(\beta)}}}{\frac{2}{\tan(\alpha)} - \frac{1}{\tan(\beta)} + \tan(\alpha)} \color{blue} \frac{\tan(\alpha)\tan(\beta)}{\tan(\alpha)\tan(\beta)} \color{black} = \frac{\tan(\alpha) \kla*{\tan(\alpha) - \tan(\beta)}}{2\tan(\beta) - \tan(\alpha) + \tan(\alpha)^2 \tan(\beta)}.
%\end{align}
\begin{align}
\tan(\epsilon) &= \tan\kla*{\arctan\kla*{\frac{1}{2a - b}} - \alpha} = \frac{\frac{1}{2a - b} - \tan(\alpha)}{1 + \tan(\alpha) \frac{1}{2a - b}} = \frac{1 - \frac{1}{a} (2a - b)}{2a - b + \frac{1}{a}} = \frac{b - a}{2a^2 - ab + 1}.
\end{align}
%\begin{tikzpicture}   
%    \pgfmathsetmacro{\a}{3.0}
%    \pgfmathsetmacro{\aa}{29}
%    \pgfmathsetmacro{\ba}{17}
%    \pgfmathsetmacro{\r}{0.05}
%    
%    
%    \coordinate (o) at (0,0);
%    \coordinate (a) at (\a,0);
%    \coordinate (a2) at ($(a) + (180-\aa:\a)$);
%    \coordinate (b) at (2*\a,0);
%    \coordinate (b2) at ($(b) + (180-\ba:\a)$);
%    \coordinate (c) at (-0.4*\a,0);
%    \coordinate (e) at (intersection of a--a2 and b--b2);
%    
%    \foreach \v in {o,a,b,e} {
%        \draw[fill] (\v) circle (\r);    
%    }
%    \draw[thick] (o) -- (a) node[below,midway] {1} -- (b) node[below,midway] {1}  -- (e) -- (o);    \draw[thick] (e) -- (a);
%    \pic["$\alpha$",draw=black,angle eccentricity=0.6,angle radius=1cm] {angle=e--a--o};
%    \pic["$\beta$",draw=black,angle eccentricity=0.8,angle radius=1cm] {angle=e--b--a};
%    \pic["$\epsilon$",draw=black,angle eccentricity=0.6,angle radius=1cm] {angle=o--e--a};
%    
%    %solution
%    \draw[thick,blue] (o) -- (c);
%    \pic["$\gamma$",draw=blue,blue,angle eccentricity=0.6,angle radius=1cm] {angle=e--o--c};
%    \node at (e) [above right,blue] {$E=(x,y)$};
%    \node at (o) [below,blue] {$O$};
%    \node at (a) [below,blue] {$A$};
%    \node at (b) [below,blue] {$B$};
%    
%    \coordinate (c2) at ($(c)!50!(o) + (\ba:\a)$);
%    \draw[thick, red] (e)+(-\aa:2cm) arc (-\aa:-42.36-\aa:2 cm);
%\end{tikzpicture}
Damit erhalten wir letztlich
\begin{align}
\epsilon &= \arctan\kla*{\frac{\cot(\beta) - \cot(\alpha)}{2\cot(\alpha)^2 - \cot(\alpha) \cot(\beta) + 1}}.
\end{align}
Für den Spezialfall $\alpha=\frac{\pi}{4}$ und $\beta=\frac{\pi}{6}$ ergibt sich damit
\begin{align}
\epsilon &= \arctan\kla*{\frac{\sqrt{3} - 1}{2 - \sqrt{3} + 1}} = \arctan\kla*{\frac{1}{\sqrt{3}}} = \frac{\pi}{6}.
\end{align}
\newpage
\section{Knifflige Winkelsumme}
Gesucht ist die Summe der drei Winkel $\alpha$, $\beta$ und $\gamma$ aus der Skizze in \autoref{figure:tikz skizze aufgabe winkelsumme}.\\
\begin{figure}[!ht]
\centering
\begin{tikzpicture}   
    \pgfmathsetmacro{\a}{3.0}
    \pgfmathsetmacro{\aa}{45}
    \pgfmathsetmacro{\ba}{30} 
    \pgfmathsetmacro{\r}{0.05}
    
    
    \coordinate (o) at (0,0);
    \coordinate (a) at (\a,0);
    \coordinate (b) at (2*\a,0);
    \coordinate (c) at (3*\a,0);
    \coordinate (e) at (0,\a);
    
    \foreach \v in {o,a,b,c,e} {
        \draw[fill] (\v) circle (\r);    
    }
    \draw[thick] (o) -- (a) node[below,midway] {1} -- (b) node[below,midway] {1} -- (c) node[below,midway] {1} -- (e) -- (o) node[left,midway] {1} ;    
    \draw[thick] (a) -- (e);
    \draw[thick] (b) -- (e);
    \pic["$\alpha$",draw=black,angle eccentricity=0.6,angle radius=1.1cm] {angle=e--a--o};
    \pic["$\beta$",draw=black,angle eccentricity=0.7,angle radius=1.2cm] {angle=e--b--o};
    \pic["$\gamma$",draw=black,angle eccentricity=0.8,angle radius=1.3cm] {angle=e--c--o};
    \pic["$\cdot$",draw=black,angle eccentricity=0.6,angle radius=0.3cm] {angle=a--o--e};
\end{tikzpicture}
\caption{Skizze der Aufgabenstellung.}
\label{figure:tikz skizze aufgabe winkelsumme}
\end{figure}
Es ist $\alpha = \frac{\pi}{4}$, da das innerste Dreieck sowohl gleichschenklig als auch rechtwinklig ist.
Schwierigkeiten bereiten nur die anderen beiden Winkel, die zunächst rein formal durch
\begin{align}
\tan(\beta) &= \frac{1}{2}\qquad \text{und} \qquad \tan(\gamma) = \frac{1}{3}
\end{align}
gegeben sind.
Tatsächlich kann man weder $\beta$ noch $\gamma$ als ein rationales Vielfaches von $\pi$ darstellen, und wir können daher nur hoffen, dass auf ``magische'' Weise deren Summe ein ``kompakteres'' Ergebnis liefert.
Es gibt hierzu eine überaus elegante Lösung mittels komplexer Zahlen.
Man beachte, dass die drei Punkte $z_1 = 1 + \i$, $z_2 = 2 + \i$ und $z_3 = 3 + \i$ gerade die Winkel $\alpha$, $\beta$ und $\gamma$ zur $x$-Achse haben.
Aufgrund der Polardarstellung komplexer Zahlen $z=r\powe{\i \phi}$ ergibt das Produkt
\begin{align}
z_1z_2z_3 &= r_1r_2r_3 \powe{\i (\alpha+\beta+\gamma)}
\end{align}
gerade eine komplexe Zahl, deren WInkel gerade der gesuchten Summe entspricht.
Im Allgemeinen hätten wir nun noch nichts gewonnen, in dem speziellen Fall ist das Produkt auf der linken Seite aber einfach
\begin{align}
z_1z_2z_3 &= (1+\i)(2+\i)(3+\i) = (1+\i) (6 + 3\i + 5\i + \i^2) = (1+\i) (5 + 5\i) = 5(1 + 2\i + \i^2) = 10\i.
\end{align}
Da dieser Wert auf der $y$-Achse liegt (oder einfach mit $\i=\powe{\i\frac{\pi}{2}}$) ist also $\alpha + \beta + \gamma = \frac{\pi}{2}$.\\
Alternativ kann man mit ``brute force'' natürlich auch auf das Ergebnis kommen.
Man beachte dass
\begin{align}
\tan(\beta + \gamma) &= \frac{\tan(\beta) + \tan(\gamma)}{1 - \tan(\beta) \tan(\gamma)} = \frac{\frac{1}{2} + \frac{1}{3}}{1 - \frac{1}{2} \frac{1}{3}} = 1
\end{align}
und somit $\beta + \gamma = \frac{\pi}{4}$.

%\begin{thebibliography}{9}
%\bibitem{example} description, \url{https://www.exmapleurl}, day.\,month~2021.
%\end{thebibliography}

\end{document}